\section{Introduction}
\label{sec:introduction}

LLMs are increasingly used as components of information and decision-support
systems that keep a task open over many turns. A procurement assistant tracks a
budget as invoices arrive; a logistics assistant follows stock across
warehouses. The user supplies information piece by piece and asks questions
whose answers depend on everything said so far, so the model's context fills
with its own earlier replies and with a task state that changes at every turn.
We call this accumulated context the \emph{conversation history}, and the
experimental condition that keeps it \emph{dialogue history}.

\paragraph{Why this is a reliability problem, not only an accuracy problem}
A growing conversation history can let a model lock onto a wrong task state:
once an early turn is answered incorrectly, every later reply can be generated
consistently from that wrong state, without the model producing any of the
usual symptoms of failure. If a deployed system does not see the gold answer,
and the responses it does see remain fluent, mutually consistent and locally
coherent, the failure is invisible unless some other signal flags it. The paper
asks two questions that follow from this chain: does retaining history cost
accuracy, and does that cost compound as the conversation grows; and if it
does, can the cost be seen before it shows up in the final answer, using
signals a deployed system can actually compute?

\paragraph{What is known}
Multi-turn use is known to be costly. Models lose accuracy when a task is
revealed over several turns because they commit to premature answers
\citep{laban2025lost}; earlier assistant replies can propagate errors
\citep{huang2026benefit}; and the effective context is shorter than the
nominal one \citep{liu2024lost,hsieh2024ruler}.

\paragraph{Two gaps these results leave open}
First, a \emph{measurement gap}: the cited studies score a single answer after
at most a handful of turns, so they show that history can hurt but not how the
loss develops over a long horizon. Over tens of turns, the questions themselves
may also become harder, so a falling accuracy curve can reflect the questions
rather than the context unless difficulty is explicitly held fixed. Second, a
\emph{monitoring gap}: none of these studies treats reasoning-quality metrics
as a stream, applies a detector calibrated on a stationary baseline, or reports
a false-alarm rate. Showing that performance eventually drops is not the same
as showing that the drop can be seen coming. ERGO is the closest exception, but
it monitors token-entropy spikes rather than response-level quality metrics,
and it is not evaluated against a long, difficulty-controlled horizon
\citep{khalid2025ergo}. To our knowledge, no prior study jointly (i) holds
difficulty fixed over a long multi-turn horizon and (ii) tests established
quality metrics as calibrated early-warning signals with a controlled
false-alarm rate.

\paragraph{This paper}
We close both gaps with two experiments, each analysed under a protocol fixed
before its outcomes were computed. Study~1 addresses the short-horizon
measurement gap directly: it compares, turn by turn, the same user information
with and without the model's own earlier replies. Study~2 addresses the
long-horizon and monitoring gaps together with CDR-Stream, a benchmark of
100-turn state-tracking dialogues with an exact answer at every turn,
difficulty calibrated per model, and a control arm that replaces the history
with a solver-verified state summary. The two studies differ in tasks, models
and controls, so we read them as two related failure modes rather than as one
dose--response curve.

The paper makes three contributions:
\begin{enumerate}
  \item a prospectively frozen, paired estimate of the cost of retaining
    assistant replies on short sharded tasks, which replicates the direction
    reported by \citet{huang2026benefit};
  \item CDR-Stream, a long-horizon, difficulty-calibrated benchmark in which
    accuracy under full history collapsed within about twenty turns at matched
    difficulty, while a verified state summary preserved it;
  \item a prospectively defined evaluation of whether established
    reasoning-quality metrics and calibrated single-stream detectors can
    serve as early warnings of this collapse -- assessing not only whether
    they move before accuracy fails, but whether any such lead survives a
    controlled false-alarm rate and repeats across models. None did: coherence
    and stability moved the wrong way, staying flat or rising as accuracy
    fell, and a closed-form bound together with a power simulation explains
    how much of this is a fundamental limit of single-conversation monitoring
    rather than a property of the metrics themselves.
\end{enumerate}
